\begin{table}[H]
\centering
\caption{Productividad laboral por departamento, 2010--2024pr}
\label{tab:pib_geih_productividad_departamento}
\scriptsize
\begin{tabular}{lrrrr}
\toprule
Departamento & PIB/trab. 2024pr & PIB/hora 2024pr & Crec. PIB/trab. & Crec. PIB/hora \\
\midrule
Quindío & 33,8 & 15,1 & 3,78\% & 4,63\% \\
Caquetá & 24,9 & 10,5 & 3,71\% & 4,17\% \\
Risaralda & 38,7 & 16,4 & 3,14\% & 3,51\% \\
Chocó & 28,6 & 11,9 & 2,95\% & 3,20\% \\
Caldas & 36,6 & 15,6 & 2,14\% & 2,57\% \\
Valle del Cauca & 54,5 & 23,4 & 2,15\% & 2,50\% \\
Sucre & 22,0 & 10,4 & 1,67\% & 2,45\% \\
Bogotá D.C. & 67,4 & 28,6 & 1,98\% & 2,44\% \\
Bolívar & 45,1 & 19,8 & 1,69\% & 2,34\% \\
Tolima & 40,0 & 16,9 & 2,04\% & 2,17\% \\
Cauca & 26,6 & 13,7 & 1,44\% & 2,05\% \\
Santander & 62,0 & 26,1 & 1,71\% & 1,96\% \\
Nariño & 17,8 & 9,1 & 0,97\% & 1,92\% \\
Córdoba & 22,0 & 10,7 & 0,90\% & 1,71\% \\
Antioquia & 47,7 & 20,3 & 1,27\% & 1,58\% \\
Boyacá & 47,5 & 22,6 & 1,16\% & 1,43\% \\
Norte de Santander & 24,3 & 10,0 & 1,23\% & 1,41\% \\
Magdalena & 23,9 & 10,3 & 0,16\% & 0,88\% \\
Atlántico & 37,3 & 16,0 & 0,46\% & 0,85\% \\
Cesar & 32,6 & 13,8 & -0,03\% & 0,57\% \\
Meta & 63,1 & 25,9 & 0,05\% & 0,55\% \\
Huila & 31,6 & 14,2 & -0,13\% & 0,11\% \\
Cundinamarca & 40,4 & 17,4 & -0,32\% & -0,14\% \\
La Guajira & 24,8 & 11,1 & -0,80\% & -0,51\% \\
\bottomrule
\end{tabular}
\caption*{\footnotesize Nota: PIB por trabajador en millones de pesos constantes de 2015; PIB por hora en miles de pesos constantes de 2015. Bogotá se trata como departamento. El universo corresponde a los 24 departamentos con información comparable en la GEIH para todo el periodo. Se excluye 2020 por no contar con GEIH anual comparable en la base del proyecto. Fuente: cálculos propios con DANE, Cuentas Nacionales Departamentales, y GEIH.}
\end{table}